\subsection{Memory Layout in Bare-Metal Compiled Processes}


\subsubsection{The Text Segment: Read-Only Machine Instructions and Instruction Pointer Tracking}

When a native executable is loaded into a process, one of the most important memory regions created by the loader is the \textbf{text segment}. Despite the name, this segment does not contain human-readable text. It contains the compiled machine instructions of the program. These are the binary instructions that the processor can fetch, decode, and execute.

In the C compilation model, the programmer writes source code such as:

\begin{verbatim}
int add(int a, int b) {
    return a + b;
}
\end{verbatim}

After compilation, assembly, and linking, this function becomes a sequence of machine instructions. Those instructions are stored in the executable file. When the operating system loader creates the process, it maps the executable instruction region into the process virtual address space. This mapped instruction region is commonly called the text segment or code segment.

A simplified process memory map may look like this:

\begin{verbatim}
high addresses
+--------------------------+
| stack                    |
+--------------------------+
| shared libraries         |
+--------------------------+
| heap                     |
+--------------------------+
| BSS                      |
+--------------------------+
| data                     |
+--------------------------+
| text / code              |
+--------------------------+
low addresses
\end{verbatim}

The exact address ranges are platform-dependent, and modern systems may randomize them for security. The conceptual division is still useful: the text segment stores executable code, while other regions store data, dynamic allocations, or stack frames.

The text segment is normally mapped as \textbf{readable and executable}, but not writable. This means the processor may read instructions from it and execute them, but ordinary program instructions should not modify it. This protection is deliberate. If program code were freely writable, then accidental memory corruption or malicious input could modify the instruction stream itself. Marking code pages as non-writable helps enforce a separation between code and data.

This permission model is part of the operating system and hardware memory-protection system. The loader maps pages with permission bits. The memory management unit checks those permissions during execution. If the program tries to write into a read-only code page, the processor traps into the kernel, and the kernel usually terminates the process with a memory protection error.

For example, a C pointer error may accidentally produce an address inside the text segment. If the program attempts to write through that pointer, the write will fail because the mapped page does not allow writing. This is not a C language rule alone. It is enforced by the process virtual memory system.

The processor executes code from the text segment by following the \textbf{instruction pointer}. The instruction pointer is a CPU register that stores the virtual address of the next instruction to execute. On x86-64, this register is usually called \texttt{RIP}. Other architectures use different names, but the role is the same.

Execution can be imagined as a repeated hardware cycle:

\begin{enumerate}
    \item read the instruction located at the address stored in the instruction pointer;
    \item decode the instruction;
    \item execute the instruction;
    \item update the instruction pointer to the next instruction, or to a different address if the instruction changes control flow.
\end{enumerate}

For straight-line code, the instruction pointer normally moves forward from one instruction to the next. For example:

\begin{verbatim}
instruction at address A
instruction at address A + n
instruction at address A + n + m
\end{verbatim}

The exact increment depends on instruction length. Some architectures have fixed-length instructions. Others, such as x86-64, have variable-length instructions. The processor still maintains the same conceptual behavior: the instruction pointer identifies the next instruction.

Control-flow instructions modify this normal movement. A conditional branch, unconditional jump, function call, or return changes the instruction pointer. For example, a C \texttt{if} statement does not remain an abstract source-level branch after compilation. It becomes comparisons and jumps. The processor follows one path or another by loading a different instruction address into the instruction pointer.

A simplified assembly shape for a condition may look like this:

\begin{verbatim}
cmp eax, 0
jle else_block
mov ebx, 1
jmp end_if

else_block:
mov ebx, -1

end_if:
\end{verbatim}

The source-level idea of an \texttt{if}/\texttt{else} structure becomes instruction-pointer redirection. If the comparison result satisfies the jump condition, the instruction pointer moves to \texttt{else\_block}. Otherwise, execution continues into the next instruction.

Function calls also operate through instruction-pointer changes. When one function calls another, the processor must jump to the callee's code while remembering where to return afterward. In many architectures, the return address is placed on the stack or in a special register. When the callee finishes, a return instruction restores the instruction pointer to the saved return address.

Thus, the text segment and the stack cooperate during function calls:

\begin{itemize}
    \item the \textbf{text segment} contains the machine instructions of both caller and callee;
    \item the \textbf{stack} stores return addresses and frame-related execution state;
    \item the \textbf{instruction pointer} moves through the text segment according to calls, branches, and returns.
\end{itemize}

For example:

\begin{verbatim}
main:
    call add
    ...

add:
    ...
    ret
\end{verbatim}

The \texttt{call} instruction changes the instruction pointer to the address of \texttt{add}. It also records the address of the instruction after the call. The \texttt{ret} instruction later uses that saved address to continue execution in \texttt{main}.

The text segment may also contain more than the programmer's direct function bodies. It can contain compiler-generated helper code, runtime startup code, procedure linkage stubs, exception-handling metadata-related code, and other low-level support routines. In a dynamically linked program, code from shared libraries is mapped into the same process address space, usually in separate text-like executable regions.

This is why a running process usually has more executable code mapped than the code written by the programmer. A small C program that calls \texttt{printf()} depends on C runtime and standard library code. A Python process contains the native code of the CPython interpreter, the C runtime, operating-system libraries, and possibly native extension modules.

The read-only nature of the text segment also allows sharing. If several processes run the same executable or use the same shared library, the operating system can map the same physical code pages into multiple process address spaces. Each process sees the code at its own virtual addresses, but the underlying physical memory containing read-only instructions may be shared. This is safe because ordinary processes are not supposed to modify those pages.

This is another reason code and data are separated. Writable data must usually be private to each process. If one process changes a global variable, that change should not automatically modify another process's data. But read-only code can be shared because it is supposed to remain identical.

The text segment therefore represents the executable core of a native process. It is the region where the compiled instruction stream lives. The processor advances through it using the instruction pointer. Branches, calls, and returns redirect that instruction pointer. The operating system protects the region as executable but normally non-writable. The result is a controlled separation between machine instructions and mutable runtime data.

This model is also useful for understanding Python. When a Python script runs, the Python source file does not become the native text segment of the process. The native text segment belongs to the CPython executable and its loaded native libraries. The Python script is compiled into Python bytecode, which is data interpreted by CPython. The physical processor executes machine instructions from CPython's text segment; CPython's interpreter loop then reads and executes Python bytecode as virtual instructions.

Thus, in C, the user's compiled functions become native instructions in the process text segment. In Python, the user's code normally becomes bytecode objects inside the interpreter runtime, while the text segment contains the interpreter's own native machine code. This distinction is one of the central memory-layout differences between bare-metal compiled execution and interpreted runtime execution.



\subsubsection{The Data and BSS Segments: Initialized and Uninitialized Global/Static Variable Storage Boundaries}

Besides the text segment, which stores machine instructions, a native compiled process also needs memory regions for values that exist for the whole lifetime of the program. In C, these are mainly global variables and static variables. They are not created and destroyed with ordinary function calls. They live for the duration of the process. Their storage is prepared when the program is loaded, before normal execution reaches \texttt{main()}.

Two traditional regions are used for this kind of storage:

\begin{itemize}
    \item the \textbf{data segment}, for initialized global and static objects;
    \item the \textbf{BSS segment}, for uninitialized or zero-initialized global and static objects.
\end{itemize}

The distinction is based on whether the executable file must physically store an initial value.

Consider the following C declarations:

\begin{verbatim}
int global_counter = 10;
static int module_flag = 1;
\end{verbatim}

Both variables have explicit initial values. Their initial bytes must be stored somewhere in the executable file, because the loader needs to place those exact values into the process memory image. These objects belong conceptually to the initialized data region.

Now consider:

\begin{verbatim}
int global_total;
static int cached_value;
\end{verbatim}

These variables have static storage duration, but no explicit initializer. In C, such objects are automatically initialized to zero before program execution begins. The executable does not need to store thousands of zero bytes for such objects. It only needs to record that a certain amount of zero-filled memory must be reserved. This region is traditionally called \textbf{BSS}.

The name BSS comes from older assembler terminology, but the modern practical meaning is simple: it is the region for static-lifetime objects that start as zero without requiring their zero bytes to be stored directly in the executable file. The loader creates the memory and ensures that it is zero-filled.

A simplified classification is:

\begin{verbatim}
int a = 5;        // initialized data
int b;            // BSS, zero-initialized

static int c = 7; // initialized data
static int d;     // BSS, zero-initialized
\end{verbatim}

This does not mean that \texttt{b} and \texttt{d} are somehow less real than \texttt{a} and \texttt{c}. Once the process is running, all of them occupy memory inside the process virtual address space. The difference is how their initial state is represented in the executable file and prepared by the loader.

The data and BSS regions are usually mapped as readable and writable. Unlike the text segment, they must be writable because the program is allowed to change global and static variables during execution. For example:

\begin{verbatim}
int counter = 0;

void increment(void) {
    counter++;
}
\end{verbatim}

The variable \texttt{counter} has static storage duration. It is not recreated every time \texttt{increment()} is called. It exists once for the whole process, and every call to \texttt{increment()} modifies the same stored object.

This is different from an automatic local variable:

\begin{verbatim}
void f(void) {
    int x = 0;
    x++;
}
\end{verbatim}

Here, \texttt{x} is normally created inside the function's stack frame each time the function is called. When the function returns, that storage is no longer part of the active function state. By contrast, a static local variable behaves differently:

\begin{verbatim}
void f(void) {
    static int x = 0;
    x++;
}
\end{verbatim}

Although \texttt{x} is only visible by name inside \texttt{f()}, its storage is not on the ordinary call stack. It has static storage duration. It is initialized once and persists across calls. The name has local visibility, but the memory lifetime is global-like.

This distinction between \textbf{scope} and \textbf{storage duration} is important. Scope answers the question:

\begin{quote}
Where in the source code can this name be used?
\end{quote}

Storage duration answers a different question:

\begin{quote}
How long does the object exist in memory?
\end{quote}

A global variable has broad scope and static storage duration. A static local variable has local scope but still has static storage duration. An ordinary local variable usually has block scope and automatic storage duration.

The data and BSS segments therefore represent memory whose lifetime is tied to the process itself, not to one function call. They are prepared before \texttt{main()} runs and remain valid until the process exits.

The executable file also benefits from separating initialized and uninitialized static data. Suppose a program declares a large global array:

\begin{verbatim}
int buffer[1000000];
\end{verbatim}

This array contains one million integers, all initially zero. If the executable had to store all those zero bytes explicitly, the binary file would become much larger. Instead, the executable records that the BSS region must reserve enough memory for the array and that this memory must begin as zero. The loader can create zero-filled memory without storing those zeros in the file.

By contrast:

\begin{verbatim}
int table[4] = {10, 20, 30, 40};
\end{verbatim}

must store the initial values \texttt{10}, \texttt{20}, \texttt{30}, and \texttt{40}. These bytes must be present in the initialized data section of the executable, because the loader needs to copy or map them into memory.

A simplified view is:

\begin{verbatim}
Executable file:
    text section      -> machine instructions
    data section      -> initial bytes for initialized globals/statics
    BSS description   -> size of zero-filled storage needed

Running process:
    text segment      -> mapped executable instructions
    data segment      -> writable initialized objects
    BSS segment       -> writable zero-initialized objects
\end{verbatim}

The data and BSS regions also interact with relocation and linking. Global variables may be referenced from different object files. A function in one translation unit may access a global variable defined in another translation unit, if it is declared with external linkage. The linker must resolve those symbols and ensure that instructions referring to them use the correct runtime addresses or relocatable access mechanisms.

For example:

\begin{verbatim}
/* globals.c */
int shared_counter = 0;
\end{verbatim}

%<<<PROGRAM:programs/mem01>>>
The object file produced from \texttt{main.c} contains a reference to \texttt{shared\_counter}. The object file produced from \texttt{globals.c} defines it. The linker connects the reference to the definition and arranges the final memory location of the variable in the data or BSS region.

The distinction between data and BSS is mostly invisible at the level of ordinary C syntax, but it matters for executable layout, memory initialization, and binary size. It explains why initialized static objects occupy space in the executable file, while large zero-initialized objects can be represented compactly.

It also helps explain why global mutable state behaves differently from stack-local state. A global or static object is not attached to the lifetime of a function call. Its address remains stable during the process lifetime, unless special relocation or loading mechanisms change the base location before execution begins. Any function that can reach the object can observe or modify the same storage.

This makes global and static mutable state powerful but dangerous. It can be used for configuration, counters, caches, shared module state, and persistent function-local state. But because the state is shared across calls and possibly across many parts of the program, it can also create hidden dependencies. A function that reads or writes a global variable may behave differently depending on previous calls elsewhere in the program.

This memory model contrasts sharply with Python's model. In Python, module-level names may look similar to C global variables, but they do not normally correspond to fixed raw storage slots in a data or BSS segment. A Python module namespace is usually a dictionary mapping names to Python objects. The objects themselves live in the CPython-managed heap. When Python code writes:

\begin{verbatim}
counter = 10
\end{verbatim}

it binds the name \texttt{counter} to an integer object. It does not reserve a fixed \texttt{int} slot in a native data segment in the same way C does.

However, CPython itself, as a native executable written largely in C, does have native data and BSS regions. The interpreter's own global C variables, runtime tables, static structures, and native library state are stored according to the native process memory model. The user's Python-level globals live above that, inside the runtime object system.

Therefore, in a bare-metal compiled C process, the data and BSS segments represent static-lifetime storage prepared by the loader before execution reaches the programmer's logic. Initialized global and static objects occupy the data segment because their initial bytes must be known. Uninitialized or zero-initialized global and static objects occupy the BSS region because only their size and zero-initialization requirement must be recorded. Together, these regions form the process-level storage boundary for values that persist independently of function call frames.


\subsubsection{The Process Heap: Dynamic Manual Memory Allocation and Deallocation Schemes (\texttt{malloc()} and \texttt{free()})}

The \textbf{heap} is the region of a process address space used for dynamic memory allocation. Unlike the text segment, data segment, BSS segment, and ordinary stack frames, the heap is not mainly determined by the static structure of the executable file. It grows and shrinks according to allocation requests made while the program is running.

In C, the heap is the memory area used when the program explicitly asks for storage whose lifetime cannot be conveniently tied to a single function call. The most familiar interface is the standard library allocation family:

\begin{verbatim}
malloc()
calloc()
realloc()
free()
\end{verbatim}

The central function is \texttt{malloc()}. It requests a block of memory of a given size and returns a pointer to the beginning of that block:

\begin{verbatim}
int *p = malloc(10 * sizeof(int));
\end{verbatim}

This line asks for enough memory to store ten integers. If the allocation succeeds, \texttt{p} receives the address of the allocated block. The program can then use that memory through pointer arithmetic or array syntax:

\begin{verbatim}
p[0] = 42;
p[1] = 99;
\end{verbatim}

This memory is not an automatic local variable. It does not disappear when the current function returns. It remains allocated until the program explicitly releases it with \texttt{free()}:

\begin{verbatim}
free(p);
\end{verbatim}

This is the core distinction between stack storage and heap storage. Stack storage is tied to call frames. Heap storage is tied to explicit allocation and deallocation.

For example:

\begin{verbatim}
void f(void) {
    int x = 10;
}
\end{verbatim}

The variable \texttt{x} is an automatic local variable. It normally lives in the stack frame of \texttt{f()}. When \texttt{f()} returns, that frame is no longer active, and \texttt{x} is no longer valid.

By contrast:

\begin{verbatim}
int *make_value(void) {
    int *p = malloc(sizeof(int));
    *p = 10;
    return p;
}
\end{verbatim}

Here, \texttt{p} itself is a local variable inside the stack frame of \texttt{make\_value()}, but the memory returned by \texttt{malloc()} is on the heap. When the function returns, the local variable \texttt{p} disappears, but the heap allocation remains. The returned pointer gives the caller a way to reach that allocated object.

This difference is essential. The pointer variable and the heap object are not the same thing. The pointer is a value that stores an address. The heap block is the memory region located at that address. The pointer may live on the stack, in global storage, inside another heap object, or in a register. The allocated block itself lives in the heap.

A simplified view is:

\begin{verbatim}
stack frame:
    p  ----+
           |
           v
heap:
    +----------------------+
    | allocated int object |
    +----------------------+
\end{verbatim}

The heap is useful because many programs need memory whose size or lifetime is not known at compile time. For example, a program may need to read a file of unknown size, build a list of unknown length, create tree nodes, allocate buffers based on user input, or create objects that outlive the function that created them.

A static array requires a known size at compile time or at least a size tied to a particular scope. A heap allocation can be decided dynamically:

\begin{verbatim}
int n;
scanf("%d", &n);

int *values = malloc(n * sizeof(int));
\end{verbatim}

Here, the size of the allocation depends on a runtime value. This is one of the major reasons heap allocation exists.

The heap is managed by an allocator. The C standard library provides the familiar interface, but the allocator itself must obtain memory from the operating system and then divide it into blocks requested by the program. On Unix-like systems, allocators may request memory from the kernel using mechanisms such as \texttt{brk()} / \texttt{sbrk()} historically, or memory mappings such as \texttt{mmap()}. On Windows, equivalent behavior is provided through Windows heap and virtual memory APIs. The exact mechanism differs by platform and allocator implementation.

From the C programmer's point of view, \texttt{malloc()} returns a usable region of memory, or returns \texttt{NULL} if the allocation fails:

\begin{verbatim}
int *p = malloc(sizeof(int));

if (p == NULL) {
    /* allocation failed */
}
\end{verbatim}

This check matters because heap memory is finite. Allocation may fail if the process cannot obtain enough memory, if address space is exhausted, or if system limits are reached.

The function \texttt{calloc()} is similar to \texttt{malloc()}, but it allocates memory for a number of elements and initializes the bytes to zero:

\begin{verbatim}
int *p = calloc(10, sizeof(int));
\end{verbatim}

The function \texttt{realloc()} changes the size of an existing allocation:

\begin{verbatim}
p = realloc(p, 20 * sizeof(int));
\end{verbatim}

This may enlarge the block in place, or it may allocate a new block, copy the old contents, free the old block, and return the address of the new one. Because the returned address may differ from the original address, \texttt{realloc()} must be used carefully.

The function \texttt{free()} releases a heap block back to the allocator:

\begin{verbatim}
free(p);
\end{verbatim}

After \texttt{free(p)}, the pointer value \texttt{p} still contains an address, but that address no longer belongs to a valid allocation owned by the program. Using it after it has been freed is an error called \textbf{use after free}:

\begin{verbatim}
free(p);
p[0] = 13;  /* invalid: use after free */
\end{verbatim}

This is one of the classic dangers of manual memory management. The compiler cannot generally prevent it, because the invalidity depends on runtime allocation history.

Another common error is a \textbf{memory leak}. A memory leak happens when allocated memory is no longer reachable by the program but has not been freed. For example:

\begin{verbatim}
void leak(void) {
    int *p = malloc(sizeof(int));
    *p = 10;
}
\end{verbatim}

When \texttt{leak()} returns, the local variable \texttt{p} disappears. If no other pointer to the allocated block exists, the program has lost the address of the heap memory. The allocator still considers the block allocated, but the program can no longer free it. Repeating this pattern can make the process consume more and more memory.

A \textbf{double free} is another serious error:

\begin{verbatim}
free(p);
free(p);  /* invalid: same allocation released twice */
\end{verbatim}

After the first call, the allocation has already been returned to the allocator. Freeing it again corrupts the allocator's understanding of heap state and may crash the program or create security vulnerabilities.

A \textbf{buffer overflow} can also occur on the heap. If the program writes beyond the allocated block, it may overwrite allocator metadata or neighboring heap objects:

\begin{verbatim}
int *p = malloc(10 * sizeof(int));
p[10] = 123;  /* invalid: valid indexes are 0 through 9 */
\end{verbatim}

The heap does not automatically know the programmer's intended array length at every access. The allocator knows the block size internally, but ordinary C pointer arithmetic does not perform automatic bounds checking. This is one of the reasons C is powerful but unsafe.

The heap allocator must track which blocks are allocated and which blocks are free. It must split large free regions into smaller blocks, merge adjacent free regions when possible, and handle fragmentation. \textbf{Fragmentation} means that free memory exists but is divided into pieces that may not satisfy a large allocation request efficiently.

There are two important forms:

\begin{itemize}
    \item \textbf{internal fragmentation}: the allocator gives a block slightly larger than requested because of alignment or size-class rules;
    \item \textbf{external fragmentation}: free memory is scattered into many separate regions.
\end{itemize}

Allocation performance also depends on allocator strategy. A simple allocator might scan a free list for a block large enough. More advanced allocators use bins, arenas, thread-local caches, size classes, or other structures to make common allocation and deallocation patterns fast. These details are implementation-specific, but the general responsibility is the same: manage variable-lifetime memory inside the process.

The heap is part of the process virtual address space. It is not a separate physical device. When the allocator needs more memory than it currently has, it asks the operating system for additional virtual memory. The kernel updates the process mappings, and physical pages may be supplied lazily when touched. Thus, even heap growth is mediated by the operating system's virtual memory subsystem.

A simplified relationship is:

\begin{verbatim}
C program
    calls malloc()
        |
        v
C library allocator
    finds or requests heap memory
        |
        v
operating system kernel
    maps additional pages if needed
        |
        v
process virtual address space
    usable heap block returned to program
\end{verbatim}

The heap also differs from the data and BSS segments. Global and static objects are allocated as part of the process image before normal execution begins. Heap objects are created during execution. Their number, size, and lifetime may depend on runtime input and control flow.

For example:

\begin{verbatim}
int global_array[100];       /* data or BSS, static lifetime */

int *dynamic_array = malloc(n * sizeof(int));  /* heap, dynamic lifetime */
\end{verbatim}

The global array exists for the entire process and has a fixed size determined by the program. The heap allocation can depend on \texttt{n}, a runtime value. It must later be released manually.

This model explains why C programmers must think carefully about ownership. If a function allocates memory, who is responsible for freeing it? The same function? The caller? Another object or subsystem? C does not enforce an ownership discipline automatically. It is a design rule imposed by the programmer.

For example:

\begin{verbatim}
char *make_buffer(size_t n) {
    char *p = malloc(n);
    return p;
}
\end{verbatim}

This function returns a heap pointer. The caller must understand that it has received ownership of an allocation that must eventually be freed:

\begin{verbatim}
char *buf = make_buffer(1024);

/* use buf */

free(buf);
\end{verbatim}

If this convention is unclear, leaks, double frees, or invalid accesses become likely.

This is a major contrast with Python. Python code also creates objects dynamically, but the programmer normally does not call \texttt{malloc()} and \texttt{free()} directly. Python objects are allocated in the CPython-managed heap. The runtime tracks references to objects, reclaims most objects through reference counting, and uses a garbage collector for certain cyclic structures.

For example:

\begin{verbatim}
a = [1, 2, 3]
\end{verbatim}

creates a Python list object and integer objects or references to integer objects inside the CPython runtime. The Python programmer does not manually free the list. When the list is no longer reachable, CPython eventually releases the memory according to its memory-management rules.

This does not mean Python has no heap. It means the heap is managed differently. In C, heap lifetime is explicit and manual. In Python, object lifetime is mediated by the interpreter's object model, reference counting, garbage collection, and allocator subsystems. The programmer still needs to understand references, aliasing, mutability, and resource lifetime, but not ordinary manual deallocation of each object.

The difference can be summarized as follows:

\begin{itemize}
    \item in C, \texttt{malloc()} creates heap storage and \texttt{free()} releases it;
    \item in Python, object creation allocates runtime-managed heap objects, and the interpreter decides when their memory can be reclaimed;
    \item in C, losing the last pointer without freeing causes a memory leak;
    \item in Python, losing the last reference usually makes the object reclaimable.
\end{itemize}

Therefore, the process heap is the memory region that supports dynamic runtime allocation. In bare-metal compiled C programs, it is controlled through explicit allocation and deallocation functions. This gives the programmer precise control over memory lifetime and layout, but also exposes the program to leaks, dangling pointers, double frees, and buffer overflows. Python preserves the idea of dynamically allocated heap objects, but moves most allocation and reclamation responsibility into the runtime system.

\subsubsection{The Process Stack: Automatic Variable Lifetime Tracking, Call Chains, and Push/Pop Stack Frames}

The \textbf{process stack} is the memory region used to track active function calls and their local execution state. In a traditional compiled C program, the stack is where the runtime records the chain of function calls currently in progress. Each active function call receives a portion of the stack called a \textbf{stack frame}. When a function is called, a new frame is created. When the function returns, that frame is removed.

This gives stack memory a very different lifetime model from heap memory. Heap memory is controlled manually in C through functions such as \texttt{malloc()} and \texttt{free()}. Stack memory is controlled automatically by the call and return structure of the program. The programmer does not normally allocate and free ordinary stack frames directly. They are created and destroyed as functions are entered and exited.

Consider a simple C program:

%<<<PROGRAM:programs/mem02>>>
When \texttt{main()} begins, a stack frame is created for \texttt{main}. That frame may contain local variables such as \texttt{a} and \texttt{b}, along with bookkeeping information required by the calling convention. When \texttt{main()} calls \texttt{square(a)}, a new stack frame is created for \texttt{square}. That frame contains the parameter \texttt{x}, the local variable \texttt{result}, and information needed to return to \texttt{main}. When \texttt{square()} returns, its frame is removed, and execution continues in \texttt{main()}.

A simplified call chain may look like this:

\begin{verbatim}
main()
    calls square()
        square() returns
main() continues
\end{verbatim}

The corresponding stack behavior is:

\begin{verbatim}
before call:

+-------------------+
| frame for main    |
+-------------------+

during square():

+-------------------+
| frame for square  |
+-------------------+
| frame for main    |
+-------------------+

after square returns:

+-------------------+
| frame for main    |
+-------------------+
\end{verbatim}

This is why the stack is often described as a \emph{last-in, first-out} structure. The most recently called function is the first one that must return. A function cannot normally return before the functions it called have returned, unless abnormal control-flow mechanisms intervene. The active call chain is therefore naturally stack-shaped.

The word \emph{push} means adding data to the top of the stack. The word \emph{pop} means removing data from the top of the stack. Function calls push new call-frame information. Function returns pop that information. The exact machine instructions and layout depend on the processor architecture, compiler, optimization level, and calling convention, but the conceptual behavior remains the same.

A stack frame may contain several categories of information:

\begin{itemize}
    \item the function's automatic local variables;
    \item copies of function arguments or references to argument locations;
    \item saved register values that must be restored after the call;
    \item the return address, meaning the instruction location where execution must continue after the function returns;
    \item bookkeeping data used for stack unwinding, debugging, or exception handling;
    \item temporary storage needed by the compiler.
\end{itemize}

The \textbf{return address} is especially important. When one function calls another, the processor must remember where to continue after the called function finishes. In many calling conventions, the return address is placed on the stack. When the called function executes a return instruction, that saved address is used to restore the instruction pointer.

For example:

\begin{verbatim}
main:
    call square
    ; execution continues here after square returns

square:
    ...
    ret
\end{verbatim}

The \texttt{call} instruction transfers control to \texttt{square} and records where to return. The \texttt{ret} instruction later restores the instruction pointer to that saved location. This is how the processor moves through the text segment while the stack records the call chain.

Two processor registers are commonly involved in stack management:

\begin{itemize}
    \item the \textbf{stack pointer}, which tracks the current top of the stack;
    \item the \textbf{frame pointer} or base pointer, which may provide a stable reference point for the current function frame.
\end{itemize}

On x86-64, these are commonly associated with registers such as \texttt{RSP} and \texttt{RBP}. Other architectures use different register names. The stack pointer changes as values are pushed or popped. The frame pointer, when used, gives the compiler a stable base from which it can address local variables and parameters using fixed offsets.

A simplified frame may look like this:

\begin{verbatim}
+----------------------------+
| previous frame information |
+----------------------------+
| return address             |
+----------------------------+
| saved frame pointer        |
+----------------------------+
| local variable 1           |
+----------------------------+
| local variable 2           |
+----------------------------+
| temporary compiler storage |
+----------------------------+
\end{verbatim}

The actual order and content are platform-dependent, but the idea is stable: a stack frame is a compact record of one active function call.

Automatic local variables in C are often stored in the function's stack frame. For example:

\begin{verbatim}
void f(void) {
    int x = 10;
    int y = 20;
}
\end{verbatim}

The variables \texttt{x} and \texttt{y} are automatic variables. Their lifetime begins when execution enters the block where they are defined and ends when that block or function returns. In many ordinary cases, their storage is part of the stack frame. When the frame is removed, the variables no longer exist as valid objects.

This is why returning the address of a local variable is invalid:

\begin{verbatim}
int *bad_function(void) {
    int x = 10;
    return &x;   /* invalid: x dies when the function returns */
}
\end{verbatim}

The pointer returned by this function points to memory that belonged to the stack frame of \texttt{bad\_function}. Once the function returns, that frame is no longer active. The address may still contain the old numeric value, but the object \texttt{x} no longer exists. Later stack activity may overwrite the same memory. Using such a pointer produces undefined behavior.

This is different from returning a heap pointer:

\begin{verbatim}
int *make_value(void) {
    int *p = malloc(sizeof(int));
    *p = 10;
    return p;
}
\end{verbatim}

Here, the local variable \texttt{p} disappears when the function returns, but the allocated heap object remains alive until \texttt{free()} is called. The lifetime of the pointer variable and the lifetime of the pointed-to object are different.

The stack is also where recursive calls accumulate. A recursive function calls itself, creating a new frame each time:

\begin{verbatim}
int factorial(int n) {
    if (n == 0) {
        return 1;
    }
    return n * factorial(n - 1);
}
\end{verbatim}

Each recursive call has its own parameter \texttt{n} and its own frame. Calling \texttt{factorial(3)} creates a chain similar to:

\begin{verbatim}
factorial(3)
    factorial(2)
        factorial(1)
            factorial(0)
\end{verbatim}

At the deepest point, several frames for the same function exist at the same time. They are separate frames, each with its own local state.

If recursion is too deep, or if functions allocate very large local objects, the stack can run out of space. This condition is usually called \textbf{stack overflow}. The stack has a limited size. Unlike the heap, it is not intended for arbitrary large dynamic storage. Large arrays or long-lived dynamically sized structures are usually better placed on the heap.

For example:

\begin{verbatim}
void dangerous(void) {
    int huge_array[100000000];
}
\end{verbatim}

This attempts to place a very large object in the function's stack frame. On many systems, it will exceed the available stack space and crash the program. A heap allocation would be more appropriate for such a large runtime object:

\begin{verbatim}
int *huge_array = malloc(100000000 * sizeof(int));
\end{verbatim}

The stack's automatic lifetime rules are useful because they make common function-local storage simple and fast. Allocating stack space is often just a matter of adjusting the stack pointer. Freeing it is often just reversing that adjustment when the function returns. This is much cheaper than general heap allocation, where the allocator must manage variable-sized blocks, free lists, fragmentation, and possible system calls.

However, the stack is less flexible than the heap. Stack objects have lifetimes tied to execution scope. They cannot safely outlive their function call unless their values are copied somewhere else. Heap objects can outlive the function that created them, but they require explicit lifetime management in C.

The process stack also plays a role in debugging and error reporting. A debugger can inspect the call stack to show which functions called which other functions. A stack trace or backtrace is a human-readable reconstruction of the active call chain. It may show, for example, that \texttt{main()} called \texttt{load\_file()}, which called \texttt{parse\_line()}, which crashed because of an invalid pointer.

A simplified backtrace might look like:

\begin{verbatim}
parse_line()
load_file()
main()
\end{verbatim}

This is possible because call-frame information and return addresses preserve the chain of active calls.

In optimized builds, compilers may change the visible stack-frame structure. They may keep some local variables in registers, omit the frame pointer, inline functions, or remove unused variables. Therefore, the conceptual model of stack frames is essential, but the exact memory layout is not guaranteed to match a simple textbook diagram in every compiled binary.

This stack model is also important when moving from C to Python. Python has function calls and call chains, but Python-level frames are not the same as raw C stack frames. CPython maintains Python frame objects or internal frame structures to represent executing Python code. These frames contain references to Python objects, local names, bytecode execution state, and links to previous frames. At the same time, the CPython interpreter itself is a native C program, so it also uses the native process stack while executing its own C functions.

Therefore, when Python code calls a Python function, two levels exist:

\begin{itemize}
    \item at the Python level, a Python execution frame is created or managed by the interpreter;
    \item at the native level, CPython's C code runs on the process stack like any other native program.
\end{itemize}

This distinction prevents a common confusion. A Python local variable is not usually a raw integer slot in the native stack like an automatic C variable. A Python local name is part of a Python execution frame and refers to a Python object stored in the managed heap. The native stack belongs to the interpreter's implementation; the Python frame belongs to the language runtime model.

Thus, the process stack is the native mechanism for tracking active function execution in compiled programs. It records call chains, return addresses, local automatic storage, saved registers, and temporary state. Its push/pop behavior matches the nested structure of function calls. It gives C efficient automatic variable lifetime management, but it also imposes lifetime limits: stack objects disappear when their frames disappear. Python preserves the idea of function frames and call chains, but it implements Python-level local variables and frames through the runtime system rather than exposing the raw native stack directly.